\documentclass[unnumsec,webpdf,modern,large]{oup-authoring-template}

% ============================================================
% Packages
% ============================================================
\usepackage[utf8]{inputenc}
\usepackage[T1]{fontenc}
\usepackage[english]{babel}
\usepackage{amsmath,amssymb,amsfonts,bm}
\usepackage{booktabs}
\usepackage{multirow}
\usepackage{float}
\usepackage{graphicx}
\graphicspath{{./}{brca_axis_figures_final/}}
\usepackage{xcolor}
\usepackage{longtable}
\usepackage{array}
\usepackage{algorithm}
\usepackage{algpseudocode}
\usepackage{tikz}
\usetikzlibrary{arrows.meta,positioning,shapes.geometric,fit,backgrounds}
\usepackage{csquotes}
\usepackage{url}
\usepackage{hyperref}
\usepackage{comment}
\setcitestyle{numbers,square,comma,sort&compress}
\bibpunct{[}{]}{,}{n}{,}{,}


\newcounter{panel}
\renewcommand{\thepanel}{\Alph{panel}}

\newcommand{\paneltag}[1]{%
    \refstepcounter{panel}%
    \label{#1}%
    \textbf{\thepanel}\par\vspace{-1.6pt}%
}

\newcommand{\panelref}[2]{Figure~\ref{#1}\ref{#2}}

\onecolumn

% ============================================================
% Compact float/table spacing for Bioinformatics
% ============================================================
\setlength{\textfloatsep}{10pt plus 2pt minus 2pt}
\setlength{\floatsep}{8pt plus 2pt minus 2pt}
\setlength{\intextsep}{8pt plus 2pt minus 2pt}
\setlength{\abovecaptionskip}{4pt}
\setlength{\belowcaptionskip}{2pt}
\renewcommand{\arraystretch}{0.93}
\setlength{\tabcolsep}{5pt}

% ============================================================
% Simple macros
% ============================================================
\newcommand{\R}{\mathbb{R}}
\newcommand{\sigmoid}{\sigma}
\newcommand{\bb}{\mathbf}
\newcommand{\softmax}{\mathrm{softmax}}

% ============================================================
% Journal metadata
% ============================================================
\journaltitle{Briefings in Bioinformatics}
\pubyear{2026}
\copyrightyear{2026}

\title[InterGATE: supervised sparse interactome learning]{InterGATE: a reproducible problem-solving protocol for supervised sparse interactome learning in transcriptomic classification}

\author[1]{Lucia Alvarez-Frutos}
\author[1,*]{Laura Senovilla}
\author[2,3,*]{Alejandro Alvaro-Meca}
\address[1]{Instituto de Biomedicina y Genética Molecular (IBGM), Consejo Superior de Investigaciones Científicas - Universidad de Valladolid, 47003 Valladolid, Spain.}
\address[2]{Centro de Investigación Biomédica en Red de Enfermedades Infecciosas (CIBERINFEC), Instituto de Salud Carlos III, Madrid, Spain.}
\address[3]{Medicina Preventiva y Salud Pública, Universidad Rey Juan Carlos, Madrid, Spain.}
\corresp[*]{Corresponding authors. Laura Senovilla, Instituto de Biomedicina y Genética Molecular (IBGM), Consejo Superior de Investigaciones Científicas - Universidad de Valladolid, 47003 Valladolid, Spain. E-mail: laurasenovilla@csic.es, Alejandro Alvaro-Meca, Medicina Preventiva y Salud Pública, Universidad Rey Juan Carlos, 28922, Alcorcón, Madrid, Spain. E-mail: alejandro.alvaro@urjc.es} 


\abstract{%
Breast cancer subtype classification from transcriptomic data requires methods that combine predictive accuracy with biological interpretability and reproducible implementation. Here, we present InterGATE, a reproducible problem-solving protocol for supervised learning of sparse, stable, and task-specific subnetworks from a prior human interactome. The workflow integrates labelled edge-type calibration, sample-conditioned edge gating, a residual multi-head graph attention backbone (ResGAT), hard TopK sparsification, stability selection, fixed-consensus fine-tuning, benchmark controls, graph export, and downstream axis-based interpretation. Applied to an integrated breast cancer transcriptomic resource of 5,812 samples and 11,907 genes, the protocol produced a 291-gene consensus graph and achieved validation and external-test macro-F1 scores of 0.928 and 0.924, respectively. The model outperformed full-expression tabular classifiers, 291-gene signature-restricted baselines, fixed-prior graph controls, rewired-prior controls, and alternative-backbone InterGATE variants under the same cohort-aware split. The implementation provides a reusable data interface for sample-by-gene expression matrices, metadata tables, and replaceable prior edge lists. Source code and scripts are available at \url{https://github.com/alex241278/InterGATE}; reproducibility data and materials are available at \url{https://doi.org/10.5281/zenodo.20815745}.
}

\keywords{graph neural networks, breast cancer, transcriptomics, interactome, network biology, reproducible workflows}


\boxedtext{%
\begin{itemize}
\item This article presents a reproducible problem-solving protocol for learning sparse, task-specific subnetworks from high-dimensional transcriptomic data and prior molecular interaction graphs.
\item InterGATE combines typed edge calibration, sample-conditioned edge gating, a residual multi-head graph attention backbone (ResGAT), hard sparsification, stability selection, benchmark controls, and graph-derived biological interpretation.
\item In a cohort-held-out external evaluation setting under unsupervised global harmonization, the proposed model outperformed full-expression tabular baselines, 291-gene signature-restricted baselines, fixed-prior graph controls, rewired-prior controls, and alternative-backbone InterGATE variants.
\item The retained 291-gene consensus graph preserves biological interpretability and the self-contained implementation can be reused with new sample-by-gene datasets and replaceable prior edge lists.
\end{itemize}}

\begin{document}
\maketitle


% ============================================================
\section{Introduction}
% ============================================================
According to the latest  GLOBOCAN estimates, there were 2.3 million new cases of breast cancer (BC) and 670,000 related deaths in 2022. Both figures are projected to increase by 38\% and 68\%, respectively, by 2050. These trends underscore the need for continued progress in diagnosis and treatment \citep{Bray2024,Freihat2025,Kim2025}.


Over the past two decades, omics technologies, particularly transcriptomics, have produced a vast amount of data. Combined with artificial intelligence (AI), these data are increasingly used to identify biomarkers that refine BC subtyping and support more precise treatment selection, advancing personalized medicine \citep{Ran2025,Sammut2022,Binder2021}. Transcriptomic studies have contributed to the identification of molecular BC subtypes, the assessment of recurrence and metastasis risk, and the prediction of treatment response. The PAM50 test analyzes the expression of 50 genes to classify BC tumors into -- luminal A (LumA), luminal B (LumB), HER2-enriched (HER2), or triple-negative (TNBC, also known as basal-like) -- while also estimating risk of recurrence (ROR) \citep{Korde2010}. MammaPrint, a 70-gene expression profile, is widely used to estimate prognosis in BC patients \citep{Glas2006}. The 21-gene Oncotype~DX assay focuses on ROR and on predicting chemotherapy benefit \citep{Cobleigh2005}. Although high-dimensional transcriptomic profiling can distinguish biologically distinct tumor states, converting these measurements into robust, interpretable models that generalize across cohorts remains a major challenge in bioinformatics \citep{Polyak2011Heterogeneity,Perou2000MolecularPortraits,TCGA2012BRCA}. 

Cancer is not solely a genetic disease. Cancer cells acquire a range of functional capabilities, including sustained proliferative signaling, evasion of growth suppressors, resistance to cell death, replicative immortality, angiogenesis, metabolic reprogramming, phenotypic plasticity, and functional heterogeneity. They also evade immune destruction and promote invasion and metastasis. These capabilities give tumors characteristic biological features such as genomic instability, non-mutational epigenetic reprogramming, tumor-promoting inflammation, vascularization, and polymorphic microbiomes. Tumor initiation and progression, however, depend not only on cancer cells but also on the surrounding microenvironment, including endothelial cells, pericytes, cancer-associated fibroblasts, macrophages, neutrophils, neurons, and senescent cells \citep{Hanahan2026}. 

Current expression-based models often ignore regulatory interactions or network topology. Earlier approaches used protein domains, gene expression, and functional annotation to predict protein-protein interactions in humans and to identify cancer-associated interaction subnetworks \citep{Rhodes2005}. However, large interactomes are not sufficient on their own. Most reference interactomes integrate evidence from multiple tissues, cell types, and experimental settings, making them context independent even though molecular interactions are highly dependent on tissue and cellular state \citep{Bossi2009,Magger2012,LamanTrip2025}. Furthermore, available interactomes are incomplete and biased toward well-studied, highly expressed proteins. This can distort topological inference and overemphasize canonical hub nodes \citep{deSilva2006,Gillis2014,Wright2024}. Many interactions also lack directionality and regulatory annotation, limiting their ability to capture causal information flow and condition-specific mechanisms \citep{Silverbush2019}. Because of these limitations, disease-oriented predictions and interpretation often require the learning or selection of task-specific subnetworks from the original interactome using cohort-derived molecular data, rather than relying on the full network.

Network medicine offers tools to study this complexity in a systematic way. It can help identify pathological nodes and pathways and reveal molecular links between apparently distinct phenotypes \citep{Barabasi2011NetworkMedicine}. Graph neural networks (GNNs) are particularly promising because they combine node-level measurements with prior molecular topology, such as protein--protein interactions and directed signaling relationships \citep{Gilmer2017MPNN,Velickovic2018GAT}. In biomedical bioinformatics, graph-based learning is increasingly used for patient stratification, biomarker discovery, and multi-omics integration, especially when prior biological networks constrain the hypothesis space \citep{Gogoshin2023GNNCancer}. Previous studies have shown the value of network-aware learning in cancer research, including network-based stratification for subtype identification \citep{Hofree2013NBS}; multi-omic graph learning with MOGONET \citep{Wang2021MOGONET}; graph-based transcriptomic prediction with post hoc subnetwork explanation in BC \citep{Chereda2021GLRP}; multimodal GNNs integrating intra-omic and inter-omic connectivity \citep{Li2024MultimodalGNN}; and large-scale masked pretraining for transcriptomic representations learning \citep{Theodoris2023Geneformer}. Even so, a unified framework that combines supervised learning of sparse subgraphs from a large prior interactome, structural stability across runs, comparative evaluation with held-out cohorts after unsupervised global harmonization, and biological interpretation of the retained graph remains largely unexplored.

Here, we study a graph learning framework that extracts a compact, reproducible, task-specific subgraph from a large human interactome. This subgraph supports multiclass discrimination, performs competitively with strong baseline classifiers, and yields a biologically structured interpretation. The framework has two layers. First, it learns a sparse supervised subgraph from a heterogeneous directed interactome. Second, it converts the retained graph into curated transcriptomic features and residual gene-enrichment scores. We frame this work as a problem-solving protocol for the recurring bioinformatics task of turning a high-dimensional expression matrix and a prior interaction network into a sparse, exportable, and biologically interpretable classification subgraph. The breast cancer analysis serves as a case study and benchmarked application of the protocol, while the accompanying software package defines the reusable input interface, execution order, configuration options, and reproducibility checks needed to apply the same workflow to other labelled transcriptomic datasets and compatible prior graphs.

% ============================================================
\section{Materials and methods}
%\section{Materials and methods}% Revisar útimos papers de bioinfortics si es system o materials
% ============================================================

\subsection{Implementation}

\subsubsection{Overview of the study design}
We used a two-layer framework to characterize multiclass BC. The first layer learned a sparse, task-specific subgraph from a large pre-existing human interactome. The second layer then interpreted the retained graph using curated transcriptomic axes and residual gene enrichment.

\begin{comment}
\subsubsection{Self-contained implementation and reusable input interface}
The implementation  was reorganized as a self-contained Python project to make the protocol executable beyond the original breast cancer case study. All paths are resolved relative to the project root, and the expected runtime layout separates processed data, raw downloaded files, external graph resources, cache files, exported artifacts, and final results. The default workflow expects two processed files, \texttt{expr\_combat\_corrected.csv} and \texttt{metadata\_combined.csv}, under \texttt{data/processed/}; however, the same software interface can be used with any sample-by-gene expression matrix provided that sample identifiers match the metadata table and the target-label column is specified in the configuration. The prior graph can also be replaced by a user-supplied edge list containing source and target genes, with optional direction, interaction type, sign, and weight fields. Dataset-specific preprocessing, such as batch harmonization, can be skipped or replaced when users provide already normalized expression matrices. Table~\ref{tab:software_interface} summarizes the input contract, configuration files, and execution entry points that define the reusable protocol interface.

\begin{table}[!htbp]
\centering
\caption{Reusable software interface for applying the protocol to the breast cancer case study or to a new labelled transcriptomic dataset. Required files define the minimum input contract; optional files make the workflow more reproducible offline.}
\label{tab:software_interface}
\resizebox{\textwidth}{!}{%
\begin{tabular}{lll}
\toprule
\textbf{Component} & \textbf{Default path or object} & \textbf{Purpose and expected values} \\
\midrule
Expression matrix & \texttt{data/processed/expr\_combat\_corrected.csv} & Samples in rows, gene symbols in columns; numeric expression values. \\
Metadata table & \texttt{data/processed/metadata\_combined.csv} & Sample identifiers, class label, cohort/batch, and optional platform fields. \\
Prior graph & \texttt{data/external/} or graph download/cache & Edge list with source gene, target gene, interaction type, optional sign/direction, and weight. \\
Configuration & \texttt{configs/default.yaml}; \texttt{intergate/config.py} & Paths, sample/label/cohort columns, execution order, cache location, and artifact location. \\
Environment & \texttt{environment.yml}; \texttt{pyproject.toml}; \texttt{requirements.txt} & Conda/pip installation, Python package metadata, and optional benchmark dependencies. \\
Data retrieval & \texttt{scripts/download\_zenodo\_data.py --extract} & Downloads Zenodo record \texttt{10.5281/zenodo.20815745}, verifies checksums when available, and arranges files. \\
Pre-flight check & \texttt{scripts/00\_check\_setup.py} & Verifies that required processed files exist and reports optional external graph resources. \\
Execution & \texttt{scripts/run\_notebooks.sh}; notebooks \texttt{0--6} & Runs graph learning, bootstrap summaries, axis analysis, residual enrichment, and benchmarks. \\
Outputs & \texttt{artifacts\_ablation/}; \texttt{Results/}; exported graph files & Stores trained bundles, consensus graph artifacts, benchmark tables, figures, and gene lists. \\
\bottomrule
\end{tabular}%
}
\end{table}
\end{comment}

\subsubsection{Datasets and preprocessing}
After harmonization and deduplication, the analysis was conducted on an integrated BC transcriptomic resource comprising 5,812 samples and 11,907 genes. The dataset combined RNA-seq and microarray cohorts and included five molecular classes: Normal (non-tumoral), LumA, LumB, HER2, and TNBC. Twelve cohorts were included: GSE1456, GSE15852, GSE162228, GSE19615, GSE20711, GSE21653, GSE25066, GSE32646, GSE58812, GSE65194, GSE96058, and TCGA. The full cohort-by-subtype breakdown is provided in Supplementary Table~1 (ST1).

Expression matrices were harmonized into a common gene symbol space. Duplicate genes were collapsed, and low-expression genes were removed using predefined filters. To reduce technical variability across studies and platforms, we applied unsupervised inter-cohort harmonization with ComBat empirical Bayes correction with batch and platform information only, without class labels as covariates \citep{Johnson2007ComBat}. Principal component analysis (PCA) and uniform manifold approximation and projection (UMAP) diagnostics showed a marked reduction in cohort- and platform-driven structure after correction, while preserving subtype-resolved biological variation (Supplementary Figure~1 (SF1) and ST2).

The final reported configuration used standard gene-wise z-score scaling, which was fitted only to the training subset:
\[
X^{(s)}_{m,i}=\frac{X'_{m,i}-\mu_i^{\mathrm{train}}}{\sigma_i^{\mathrm{train}}}.
\]
 The resulting matrix \(X^{(s)}\) was used as the input at the node level. A workflow overview, training pseudocode, and main hyperparameters are provided in SF2, Supplementary Algorithm~1 (SA1), and ST3.

\subsubsection{Machine-learning evaluation design, external test, and preprocessing control}
To prevent information leakage between development and evaluation cohorts, supervised analyses used cohort-aware partitioning for the external test set. Cohort, rather than individual samples, defined the final external split. The integrated resource was first divided, using a fixed random seed, into a non-test pool (\emph{train\_pool}, \(n=4{,}757\)) and a held-out external test partition (\(n=1{,}055\)). The external test set comprised three cohorts: TCGA (\(n=904\)), GSE58812 (\(n=107\)), and GSE32646 (\(n=44\)). A validation subset was then sampled from the non-test \emph{train\_pool}, yielding 3,805 training samples and 952 validation samples. Validation was used for early stopping, model selection, and baseline tuning, whereas the external test set was reserved exclusively for final assessment. Because inter-cohort harmonization was performed globally before partitioning, the evaluation should be interpreted as cohort-held-out under unsupervised global harmonization rather than as strictly leakage-free pipeline. Gene-wise z-score scaling was fitted to the training data and applied unchanged to the validation and test data. 

\subsubsection{Prior graph construction}
A directed prior graph was built by integrating two complementary human interactome resources: HuRI \citep{Luck2020HuRI}, which mainly provides undirected physical protein--protein interactions, and OmniPath \citep{Turei2016OmniPath}, which provides directed signaling interactions with optional activating/inhibitory annotations. After restriction to the study-specific gene universe, HuRI contributed 35,118 directed edges after bidirectional expansion, and OmniPath contributed 17,186 directed interactions. After deduplication, the final integrated prior contained 51,445 directed edges. Removing isolated nodes yielded a connected backbone retaining 7,213 genes, corresponding to 60.6\% of the 11,907 genes in the expression universe. The remaining 4,694 genes (39.4\%) were excluded from graph-based message passing. The final backbone matrix used for graph learning therefore had dimensions \(5,812 \times 7,213\). Self-loops were not added, and trivial self-edges were excluded from the pruning budget.

\subsubsection{Sample-level aggregated covariates}
In addition to gene-level expression, we defined a sample-level covariate matrix \(X_h\), to provide global context. For each regulator in the prior graph, outgoing targets were collected, and only regulators with at least five targets were retained. For each selected regulator and each sample, summary statistics (mean, standard deviation, and maximum) were computed over the scaled expression of its targets, yielding:
\[
X_h \in \mathbb{R}^{M\times F}.
\]
These covariates were used to condition edge gates during graph learning, which allowed for sample-specific modulation of edge scores, and to enrich the graph-level classifier during fine-tuning, if desired. To avoid circularity, \(X_h\) was computed once from the fixed prior graph and never recomputed after subgraph selection. 

\subsection{Algorithm}

\subsubsection{Formal definition of the ResGAT backbone with typed edge gating}
Let \(M\) denote the number of samples, \(C\) the number of classes, and \(G_0=(V,E)\) the directed candidate interactome after restriction to the study-specific gene universe. Each node \(v\in V\) represents one gene and each directed edge \(e=(u,v)\in E\) represents a candidate interaction from regulator \(u\) to target \(v\). Each edge has a discrete interaction type \(t_e\), a fixed prior weight \(w_e^{(0)}\), and, when available, a sign distinguishing activating and inhibitory evidence.

For sample \(m\), the scaled expression of gene \(v\) is denoted by \(x^{(s)}_{m,v}\). This scalar input is projected to an initial latent embedding through a shared input multilayer perceptron:
\[
h^{(0)}_{m,v}=\phi_{\mathrm{in}}\!\left(x^{(s)}_{m,v}\right)\in\mathbb{R}^{d}.
\]
The graph backbone is a residual multi-head graph attention network (ResGAT) implemented in a general message-passing framework \citep{Gilmer2017MPNN,Velickovic2018GAT}. Messages are aggregated over the incoming neighborhood of each target node in the directed graph.

To learn a task-specific subgraph, each candidate edge \(e\) is assigned a trainable global gate logit \(\theta_e\). Because the prior graph combines multiple edge types, this logit is first calibrated by a type-specific affine transformation,
\[
\theta_e^{\mathrm{type}}=a_{t_e}\theta_e+b_{t_e},
\]
where \(a_t\) and \(b_t\) are learned separately for each interaction type. This allows the model to calibrate the ranking scale differently for physical interactions, activating edges, and inhibitory edges.

To model inter-sample heterogeneity, we compute a sample-dependent offset from the regulator-summary covariates \(X_h\):
\[
\Delta(X_{h,m})=(\delta_{m,0},\delta_{m,1},\delta_{m,2}),
\]
where \(\Delta(\cdot)\) is a lightweight MLP. The sample-specific pre-sigmoid score for edge \(e\) in sample \(m\) is then
\[
s_{m,e}(\tau)=\frac{\theta_e^{\mathrm{type}}}{\tau}+\delta_{m,t_e},
\]
where \(\tau\) is the temperature. The corresponding soft gate is
\[
g_{m,e}(\tau)=\sigma\!\left(s_{m,e}(\tau)\right),
\]
and the effective sample-specific edge weight becomes
\[
\bar w_{m,e}(\tau)=w_e^{(0)}\,g_{m,e}(\tau).
\]

During supervised graph learning, hard structural sparsity is enforced by retaining only a fraction \(\rho\) of edges with the highest global calibrated pre-sigmoid scores. At each forward pass, a global TopK mask is computed using \(\theta_e^{\mathrm{type}}/\tau\), producing a retained edge set \(E_{\mathrm{keep}}\subseteq E\). This hard pruning step is shared by all samples in the batch, whereas the sample-level offsets modulate only the soft gates of already retained edges.

To stabilize propagation, retained edge weights are symmetrically normalized using the weighted magnitudes of incident edges. For node \(u\), the sample-specific outgoing soft degree is
\[
\hat d^{\mathrm{out}}_{m,u}(\tau)=\sum_{e=(u,j)\in E_{\mathrm{keep}}}\left|\bar w_{m,e}(\tau)\right|,
\]
and for node \(v\), the incoming soft degree is
\[
\hat d^{\mathrm{in}}_{m,v}(\tau)=\sum_{e=(i,v)\in E_{\mathrm{keep}}}\left|\bar w_{m,e}(\tau)\right|.
\]
The normalized effective edge weight used during message passing is
\[
\omega_{m,(u,v)}(\tau)=
\frac{\bar w_{m,(u,v)}(\tau)}{\sqrt{\left(\epsilon+\hat d^{\mathrm{out}}_{m,u}(\tau)\right)\left(\epsilon+\hat d^{\mathrm{in}}_{m,v}(\tau)\right)}},
\]
where \(\epsilon>0\) is a small numerical constant. Absolute magnitudes are used in the degree terms so that activating and inhibitory contributions do not cancel each other during normalization.

For ResGAT layer \(\ell\) and attention head \(k\), node embeddings are linearly transformed as
\[
z^{(\ell,k)}_{m,v}=W^{(\ell,k)}h^{(\ell)}_{m,v}.
\]
For each retained edge \((u,v)\in E_{\mathrm{keep}}\), the unnormalized attention score is
\[
a^{(\ell,k)}_{m,(u,v)}=
\mathrm{LeakyReLU}\!\left[
\mathbf{q}^{(\ell,k)\top}
\Big(z^{(\ell,k)}_{m,u}\,\Vert\,z^{(\ell,k)}_{m,v}\,\Vert\,\omega_{m,(u,v)}(\tau)\Big)
\right],
\]
where \(\Vert\) denotes concatenation. Attention coefficients are normalized over the incoming neighbors of node \(v\):
\[
\alpha^{(\ell,k)}_{m,(u,v)}=
\frac{\exp\left(a^{(\ell,k)}_{m,(u,v)}\right)}{\sum_{r:(r,v)\in E_{\mathrm{keep}}}\exp\left(a^{(\ell,k)}_{m,(r,v)}\right)}.
\]
The head-specific aggregated message received by node \(v\) is
\[
\mu^{(\ell,k)}_{m,v}=
\sum_{u:(u,v)\in E_{\mathrm{keep}}}
\alpha^{(\ell,k)}_{m,(u,v)}\,\omega_{m,(u,v)}(\tau)\,z^{(\ell,k)}_{m,u}.
\]
Finally, the outputs of the \(K\) heads are concatenated and passed through a residual block:
\[
h^{(\ell+1)}_{m,v}=\mathrm{LayerNorm}\!\left[
R^{(\ell)}h^{(\ell)}_{m,v}+\mathrm{Dropout}\!\left(
\psi^{(\ell)}\!\left(\Vert_{k=1}^{K}\mu^{(\ell,k)}_{m,v}\right)
\right)
\right],
\]
where \(R^{(\ell)}\) is either the identity map or a learned linear projection when dimensions differ, and \(\psi^{(\ell)}\) is a feed-forward transformation. This formulation defines the ResGAT backbone used in the primary InterGATE configuration.

\subsubsection{Readout and hybrid classifier}
After \(L\) ResGAT layers, node embeddings were aggregated into a graph-level representation using multi-head global attention pooling. In the graph-only branch, this embedding could be concatenated with \(X_h\) and passed through an MLP classifier. Once the consensus subgraph was established, we activated a hybrid classification head combining: (i) the graph branch (ResGAT + pooling); (ii) a tabular MLP branch that operates directly on the compact expression vector; and (iii) a fusion block. In the primary configuration, the latent vectors from both branches were concatenated and processed by a final fusion MLP to produce class logits. 

\subsubsection{Training procedure}
\label{subsec:training_procedure}
Training was carried out in four stages. First, an optional self-supervised pretraining module based on masked gene modeling was implemented, in line with masked reconstruction methods in graph self-supervision and transcriptomic foundation modeling \citep{Hou2022GraphMAE,Theodoris2023Geneformer}; however, pretraining was disabled in the primary configuration. Second, supervised graph learning was run for up to 90 epochs, with temperature \(\tau\) and the retained-edge fraction \(\rho\) following predefined schedules. \(\rho\) decreased exponentially from 1.0 to the final sparsity floor \(\kappa=0.004\). Optimization used AdamW \citep{Kingma2015Adam,Loshchilov2019AdamW} with gradient accumulation, and validation performance was monitored using macro-F1 and early stopping. During this phase, \(X_h\) was disabled in the classifier so that structural selection depended mainly on gene expression, although it could still condition edge gates. Third, stability selection was performed over \(R=8\) independent runs. In each run, the model was retrained with fixed hyperparameters for 60 epochs, and the final graph containing the top \(\kappa=0.004\) fraction of edges by calibrated pre-sigmoid ranking was exported. Edge selection frequency \(\pi_e\) was then computed across runs. The consensus graph retained edges with \(\pi_e \ge 0.50\), which is the threshold recommended in the original stability selection framework \citep{Meinshausen2010Stability}. The graph was then compacted by removing genes not incident to any retained edge. Fourth, the model was fine-tuned on the fixed consensus graph in two stages: an initial 5-epoch stage without the hybrid branch and a longer 80-epoch stage with the hybrid GNN+MLP head and optional mixup augmentation \citep{Zhang2018Mixup}. During mixup, interpolation was applied only to sample features, while keeping the graph structure fixed.

All experiments were run on a single workstation with an Intel Core i9-14900K processor, 192 GB RAM, and an NVIDIA GeForce RTX 3090 GPU. The full stability-selection pipeline, including \(R=8\) independent supervised graph-learning runs, took approximately 200 minutes.


\subsubsection{Losses, model selection, and evaluation} 
The main supervised objective was weighted multiclass cross-entropy, which accounts for class imbalance:
\[
\mathcal{L}_{\mathrm{sup}}(m)= -\sum_{c=0}^{C-1}\omega_c\,\mathbf{1}[y_m=c]\log \hat p_{m,c}.
\]

During phase 1 graph learning, this objective was optimized alongside a sparsity regularizer on the global edge gates. When enabled, a connectivity penalty was also applied to reduce topological collapse during aggressive pruning:
\[
\mathcal{L}_{\mathrm{phase1}}
=
\mathcal{L}_{\mathrm{sup}}
+
\lambda_{\mathrm{gate}}\mathcal{R}_{\mathrm{gate}}(\tau)
+
\lambda_{\mathrm{conn}}\mathcal{R}_{\mathrm{conn}}(\tau),
\]
with
\[
\mathcal{R}_{\mathrm{gate}}(\tau)=\sum_{e\in E}\sigma\!\left(\frac{\theta_e^{\mathrm{type}}}{\tau}\right).
\]
The connectivity term was defined on the directed candidate graph \(G=(V,E)\), where \(v\in V\) denotes a gene node and \(e=(u,v)\in E\) denotes a directed candidate interaction. Using the global, non-sample-conditioned gate
\[
g_e(\tau)=\sigma\!\left(\frac{\theta_e^{\mathrm{type}}}{\tau}\right),
\]
we first computed a soft magnitude for each edge,
\[
\tilde w_e(\tau)=\left|w_e^{(0)}\right|g_e(\tau),
\]
where \(w_e^{(0)}\) is the fixed prior edge weight. The expected outgoing and incoming soft degrees of node \(v\) were then
\[
\hat d_v^{\mathrm{out}}(\tau)=\sum_{e=(v,u)\in E}\tilde w_e(\tau),
\qquad
\hat d_v^{\mathrm{in}}(\tau)=\sum_{e=(u,v)\in E}\tilde w_e(\tau),
\]
and the total incident soft degree was
\[
\hat d_v(\tau)=\hat d_v^{\mathrm{out}}(\tau)+\hat d_v^{\mathrm{in}}(\tau).
\]
The connectivity penalty was therefore
\[
\mathcal{R}_{\mathrm{conn}}(\tau)=
\frac{1}{|V|}\sum_{v\in V}\left[d_{\min}-\hat d_v(\tau)\right]_+,
\]
where \([x]_+=\max(x,0)\) and \(d_{\min}\) is the minimum expected soft-degree threshold. Thus, nodes whose incident gated edge mass fell below \(d_{\min}\) were penalized, whereas sufficiently connected nodes contributed zero to this term. In the reported configuration, the degree was computed using absolute prior edge weights so that activating and inhibitory signed interactions did not cancel each other.

After consensus-graph fixing, phase 2 fine-tuning used weighted cross-entropy on the fixed graph. Mixup was applied only during the final sub-stage fine-tuning when enabled. The gate sparsity term was not used after graph selection was complete. Model selection in phase 1 relied on validation macro-F1 with early stopping. In phase 2, the best checkpoint was selected again based on validation macro-F1. Final performance was evaluated once on a held-out external test set, which was not used during supervised training or model selection, using accuracy, macro-F1, weighted F1, and, when available, OvR multiclass AUC. 

\subsubsection{Comparative benchmark protocol against reference baselines}
To benchmark the method, we compared it with several model families using the same cohort-aware training, validation, and testing split; preprocessing protocol; and label space. First, we trained full-expression tabular baselines on the harmonized matrix comprising all 11,907 genes, including multinomial elastic-net logistic regression, calibrated linear support vector machine (SVM), graph-free multilayer perceptron, random forest, extremely randomized trees, \texttt{XGBoost}, \texttt{LightGBM}, and \texttt{CatBoost} \citep{Ke2017LightGBM,Prokhorenkova2018CatBoost}. Hyperparameters were selected using only the validation set. Second, to distinguish the impact of topology-aware graph learning from simple compact feature restriction, we trained additional tabular baselines restricted to the final 291-gene consensus signature. The signature-restricted models included a random forest, \texttt{XGBoost}, elastic-net logistic regression, and a graph-free multilayer perceptron. These models were selected by validation macro-F1. Third, we trained fixed-graph controls based on the same prior-graph family but with graph learning disabled. In these controls, we kept the prior graph fixed, neutralized and froze edge gating, and optimized only the downstream classifier parameters. We considered graph-only and hybrid graph+tabular variants. Fourth, we evaluated rewired-prior controls to determine if performance could be explained by an arbitrary graph of the same size and composition. The prior graph was randomized while preserving the node set, edge count, edge-type counts, and edge weight/sign assignments. The resulting rewired graph was then used as a fixed backbone in the graph-only and hybrid variants. Fifth, all benchmark models---tabular, signature-restricted, fixed-prior, and rewired-prior---were selected by validation macro-F1 and evaluated on the same held-out external test set. This expanded design assesses whether supervised learning of a sparse, task-specific interactome subgraph provides added value beyond strong full-expression tabular models, compact signature-restricted tabular models, fixed biological graph controls, and arbitrary graph controls.


\subsubsection{Comparative ablation and selection of the primary reported configuration}
We also evaluated a predefined set of internal architectural variants to measure the impact of critical design decisions, such as edge-type calibration, sample-conditioned gating, signed modeling, and connectivity regularization. To ensure that performance differences could be attributed to the architectural change itself, all variants were run under the same data partition, optimization schedule, and random seed. Comparisons were performed on the validation set, while the external test set was reserved for the final evaluation of the selected model and was not used for configuration selection. The primary reported configuration was selected according to validation macro-F1. In addition, we conducted a single-seed backbone-replacement diagnostic in which the FULL protocol and hybrid head were kept active while replacing only the message-passing backbone with GraphSAGE \citep{Hamilton2017GraphSAGE}, GIN \citep{Xu2019GIN}, or GraphTransformer-style message passing \citep{Shi2020TransformerConv,Dwivedi2020GraphTransformer}. This analysis was treated as an internal architectural robustness check and is reported in the supplementary material, separately from the main benchmark panel.

\subsubsection{Definition of graph-derived refined transcriptomic axes}
A second downstream layer was defined using curated transcriptomic axes derived from the retained graph signature. The final, validation-selected, consensus graph contained 291 genes and the full list is provided in Supplementary Data File~1 (SD1). Based on these genes, we manually curated eight refined biological axes:
\begin{enumerate}
    \item \texttt{Luminal Hormonal}
    \item \texttt{Cell Cycle Mitotic}
    \item \texttt{HER2 RTK MAPK}
    \item \texttt{Basal Plasticity TNBC}
    \item \texttt{Immune Lymphoid Signaling}
    \item \texttt{DNA Damage p53 Checkpoint}
    \item \texttt{Adhesion Cytoskeleton Invasion}
    \item \texttt{Androgen Apocrine}
\end{enumerate}

Rather than treating these axes as external canonical signatures, they were treated as curated, graph-constrained interpretive modules. Each axis grouped genes associated with recurrent biological programs reported in the breast cancer literature. Axes were permitted to overlap when a retained gene plausibly connected multiple programs. Specifically, \texttt{EGFR} appears in both the HER2 RTK MAPK and Basal Plasticity TNBC axes, and \texttt{MSLN} appears in both the Basal Plasticity TNBC axes and Androgen Apocrine axes. The complete gene composition and axis sizes are provided in ST4. This design ensures that the downstream interpretation aligned with the learned graph structure rather than with externally defined gene panels.

\subsubsection{Axis coverage, standardization, and sample-level scoring}
Axis coverage was evaluated with both the retained graph signature, which indicates which biological programs are structurally present in the learned subgraph, and the full expression matrix, which ensures that axis scores can be computed across all samples. This axis-level standardization was only used for the downstream interpretability layer and was not part of the training, validation, or external test evaluation of the predictive model. For each gene \(g\), expression was standardized across samples using z-score normalization:
\[
z_{ig} = \frac{x_{ig}-\mu_g}{\sigma_g},
\]
where \(x_{ig}\) denotes the expression of gene \(g\) in sample \(i\). For each biological axis \(k\), defined by gene set \(G_k\), a per-sample axis score was computed as the average standardized expression across all genes in the set:
\[
s_i^{(k)} = \frac{1}{|G_k|}\sum_{g\in G_k} z_{ig}.
\]

\subsubsection{Statistical analysis of refined axes}
To assess subtype-associated differences in axis activity, we employed a non-parametric, two-level statistical approach. First, we performed a global Kruskal--Wallis test across the five molecular subtypes for each axis. Second, we performed an OvR analysis and a two-sided Mann--Whitney U test for each subtype (\texttt{Normal}, \texttt{Luminal A}, \texttt{Luminal B}, \texttt{HER2-enriched}, \texttt{TNBC}). In addition to raw \(P\)-values, we reported effect size using the mean difference, rank-biserial correlation, and Cohen's \(d\). 

\subsubsection{Residual gene set definition and functional enrichment} 
After defining the eight axes, the 291-gene signature was divided into genes assigned to the curated axes and residual genes not assigned to any axis. The final enrichment analyses considered 83 assignments corresponding to 81 unique genes, leaving 210 residual genes. Functional enrichment of these residual genes was performed with GSEApy through the Enrichr interface \citep{Fang2023GSEApy,Kuleshov2016Enrichr}, using HGNC gene symbols as input. The queried libraries included \texttt{KEGG 2021 Human}, \texttt{Reactome 2022}, \texttt{GO Biological Process 2023}, \texttt{GO Molecular Function 2023}, and \texttt{GO Cellular Component 2023}.

% ============================================================
\section{Results}
% ============================================================
\subsection{Extended benchmarking against tabular, signature-restricted, and graph-structure controls}
Using the same cohort-aware split and preprocessing pipeline, the proposed graph-learning model remained the best in the extended benchmark panel (Table~\ref{tab:benchmark_panel_extended}). The validation-selected \texttt{FULL} configuration achieved validation and external test macro-F1 scores of 0.928 and 0.924, respectively with an accuracy score of 0.948 and an OvR macro-AUC score of 0.991 (more details in ST5).

Among the full-expression tabular baselines, \texttt{random\_forest} was the strongest model on the external test set (macro-F1 score, 0.902), followed by \texttt{CatBoost} and \texttt{LightGBM} (both 0.896) and \texttt{extratrees} (0.894). In contrast, \texttt{XGBoost} achieved the highest validation macro-F1 score in this group (0.975), but decreased to 0.837 on the external test set. We trained tabular baselines on the final 291-gene consensus set to determine if the improvement could be attributed solely to compact feature restriction. In this setting, \texttt{random\_forest\_291} reached an external test macro-F1 score of 0.900, while \texttt{XGBoost\_291} reached 0.893. Both remained below the graph-learning model's performance.

Fixed-prior controls from the same prior-graph family performed substantially worse than the proposed model, achieving an external test macro-F1 score of 0.827. Rewired-prior controls also remained below the proposed method, reaching external test macro-F1 scores between 0.843 and 0.846. These results suggest that the observed improvement cannot be attributed solely to high-capacity tabular classification, compact feature restriction, a fixed biological prior, or arbitrary graph structure. Rather, the results are most consistent with a contribution from supervised, topology-aware learning of a sparse, task-specific subgraph.

\begin{table}[!htbp]
\centering
\caption{Extended benchmark comparison under the same cohort-aware split. Representative baselines are shown from three families: full-expression tabular, 291-gene signature-restricted tabular, and graph-based controls. The graph-control panel includes fixed-prior and rewired-prior controls. Alternative message-passing backbones within the FULL InterGATE protocol are reported separately as an internal architectural diagnostic in ST6.}
\label{tab:benchmark_panel_extended}
\resizebox{\textwidth}{!}{%
\begin{tabular}{llcccc}
\toprule
\textbf{Family} & \textbf{Model} & \textbf{Validation macro-F1} & \textbf{Test accuracy} & \textbf{Test macro-F1} & \textbf{Test OvR macro-AUC} \\
\midrule
Proposed graph-learning
& \texttt{FULL} & 0.928 & 0.948 & \textbf{0.924} & \textbf{0.991} \\
\midrule
\multirow{5}{*}{Full-expression tabular}
& \texttt{random\_forest}        & 0.938 & 0.918 & 0.902 & 0.983 \\
& \texttt{CatBoost}              & 0.964 & 0.892 & 0.896 & 0.984 \\
& \texttt{extratrees}            & 0.860 & 0.910 & 0.894 & 0.986 \\
& \texttt{elasticnet\_logreg}    & 0.862 & 0.889 & 0.870 & 0.987 \\
& \texttt{XGBoost}               & 0.975 & 0.848 & 0.837 & 0.971 \\
\midrule
\multirow{3}{*}{291-gene signature tabular}
& \texttt{random\_forest\_291}     & 0.943 & 0.942 & 0.900 & 0.990 \\
& \texttt{XGBoost\_291}            & 0.953 & 0.912 & 0.893 & 0.985 \\
& \texttt{elasticnet\_logreg\_291} & 0.787 & 0.865 & 0.844 & 0.981 \\
\midrule
\multirow{3}{*}{Graph controls}
& \texttt{fixed\_prior\_hybrid}       & 0.845 & 0.848 & 0.827 & 0.972 \\
& \texttt{rewired\_prior\_gnn}        & 0.770 & 0.867 & 0.846 & 0.984 \\
& \texttt{rewired\_prior\_hybrid}     & 0.851 & 0.864 & 0.843 & 0.973 \\
\bottomrule
\end{tabular}%
}
\end{table}


\subsection{Comparative ablation on the validation set selected the primary reported configuration}
We compared seven predefined model variants with the same seed, data split, and training schedule (Table~\ref{tab:ablation_configs}). We assessed performance on the validation set.

Of these variants, \texttt{FULL} showed the strongest overall performance, with validation accuracy of 0.953, a macro-F1 of 0.928, and OvR macro-AUC of 0.994. Therefore, we used this validation-selected configuration to derive the final 291-gene consensus graph for downstream structural and biological analyses.

The remaining variants performed less well. This suggests that using typed edge calibration, sample-conditioned gating, signed modeling, and connectivity regularization together had a positive effect in the primary reported setting.

\begin{table}[!htbp]
\centering
\caption{Comparative ablation across seven predefined architectural variants. The reported metrics correspond to validation-set performance and were used to select the primary configuration.}
\label{tab:ablation_configs}
\resizebox{\textwidth}{!}{%
\begin{tabular}{lcccc}
\toprule
\textbf{Configuration} & \textbf{Validation accuracy} & \textbf{Validation macro-F1} & \textbf{Validation OvR macro-AUC} & \textbf{Final signature genes} \\
\midrule
\texttt{FULL} & 0.953 & 0.928 & 0.994 & 291 \\
\texttt{no\_signed} & 0.909 & 0.885 & 0.987 & 317 \\
\texttt{dual\_backbone} & 0.905 & 0.883 & 0.986 & 295 \\
\texttt{no\_conn} & 0.881 & 0.853 & 0.979 & 301 \\
\texttt{typed\_only} & 0.860 & 0.840 & 0.978 & 299 \\
\texttt{no\_type\_no\_cond} & 0.841 & 0.821 & 0.974 & 332 \\
\texttt{cond\_only} & 0.802 & 0.795 & 0.971 & 283 \\
\bottomrule
\end{tabular}
}
\end{table}

The notably lower performance of \texttt{cond\_only} further suggests that sample-conditioned modulation alone is insufficient for robust graph selection. Without typed calibration, instance-level offsets do not seem to compensate for the lack of a stable global edge-ranking mechanism.

A separate backbone-replacement diagnostic is reported in ST6. Keeping the FULL protocol and hybrid head active, the GAT/ResGAT variant achieved the highest external test macro-F1 score (0.924), followed by GraphTransformer (0.841), GIN (0.834), and GraphSAGE (0.830). Because this diagnostic used one seed and did not replace the full stability-selected consensus procedure, it supports the choice of ResGAT but was not used to redefine the primary 291-gene model.

\subsection{The selected primary configuration showed favorable multiclass classification performance on the external test set}
Classification performance was evaluated using the validation-selected primary configuration (\texttt{FULL}) and its final 291-gene consensus graph, on the external test set (\(n=1{,}055\) samples). This set comprised three held-out cohorts (TCGA, GSE58812, and GSE32646), which were kept separate from all supervised model development steps, including training, model selection, and comparative ablation.

On the external test set, the model achieved an accuracy of 0.948, a macro-precision of 0.929, a macro-recall of 0.920, and a macro-F1 score of 0.924. The weighted F1 score was 0.947, and the OvR macro-AUC was 0.991. Detailed, class-specific metrics and bootstrap confidence intervals are provided in ST7. The observed confusion matrix and bootstrap OvR AUC estimates are provided in ST8 and ST9, respectively. The model performed best on the luminal classes. HER2-enriched and healthy tissue showed lower F1 values and wider uncertainty intervals, consistent with their smaller sample size and partial overlap with neighboring groups.

The external test macro-F1 score was slightly lower than the validation macro-F1 score (0.924 versus 0.928), with a difference of 0.004. In a single-cohort held-out evaluation, this gap is small and consistent with typical variability in cohort composition and class difficulty. Since the external test partition remained unchanged throughout model development and was only used once for the final assessment, this pattern suggests generalization to held-out cohorts rather than selection bias.

\subsection{The final retained graph defined a sparse, task-specific regulatory network}
The prior graph was constructed by combining HuRI and OmniPath across \(11{,}907\) genes. This yielded \(7{,}213\) connected nodes (\(60.6\%\)) and \(51{,}445\) directed edges after deduplication. Following supervised graph learning and stability selection across \(R=8\) independent runs, the consensus subgraph retained 291 genes. This represents a significant compression of the original interactome while preserving a biologically interpretable regulatory backbone for discriminating between multiclass subtypes.

The final retained graph is shown in Figure~\ref{fig:final_graph}. The network is organized around a central, signaling-rich core with prominent hubs, including \texttt{ESR1}, \texttt{EGFR}, \texttt{ERBB2}, \texttt{TP53}, and \texttt{ABL1}. This core is surrounded by a broader periphery of lower-degree nodes. Each of the 8 stability run exports contained 205 edges, and the overlap with the final consensus graph ranged from 31.6\% to 42.1\% (mean 37.8\%), indicating the recovery of a reproducible core subgraph.

\begin{figure}[!ht]
    \centering
    \includegraphics[width=0.9\textwidth]{Figure1.png}
    \caption{Final retained regulatory graph after supervised graph learning and stability selection across \(R=8\) runs. The size of the nodes in the retained graph is proportional to their degree. Darker colours correspond to smaller nodes and lighter colours to larger ones. The most connected hubs, including \texttt{ESR1}, \texttt{EGFR}, \texttt{ERBB2}, \texttt{TP53}, and \texttt{ABL1}, occupy the central core of the network. Edge colour encodes interaction type: blue for HuRI protein--protein interactions, green for OmniPath activation edges, and red for OmniPath inhibition edges.}
    \label{fig:final_graph}
\end{figure}

\newpage
\subsection{Axis coverage and subtype-level structure remained biologically coherent}
All eight curated axes were represented in the retained 291-gene graph and remained directly measurable in the full expression matrix. Therefore, the downstream axis-based layer relied on modules anchored in the learned signature itself rather than on partial projection onto external signatures.

Mean axis scores by subtype are summarized in ST10. These values revealed a clear, biologically consistent pattern. Luminal A was characterized by strong activation of the luminal axis and repression of the mitotic axis. Interestingly, the Luminal B subtype was less strongly characterized by activation of the luminal axis. Its profile shifted toward a more proliferative state. Both luminal subtypes showed suppression of the basal-TNBC axes. The HER2-enriched subtype showed reduced luminal signaling, activated proliferation, and partial activation of the immune signaling pathway. Finally, TNBC exhibited the most extreme pattern. It combined strong repression of the luminal axis with activation of the cell cycle, basal-TNBC, and immune signaling axes. TNBC also showed partial activation of the adhesion-invasion program (SF3). 

\subsection{Integrated visual readout of subtype geometry}
The integrated view of the subtype geometry includes clustering at the subtype level (Figure~\ref{fig:axis_heat_corr_pca}A), consistent relationships among the refined axes (Figure~\ref{fig:axis_heat_corr_pca}B), partial overlap between biologically adjacent groups (Figure~\ref{fig:axis_heat_corr_pca}C), and differences between the subtypes across the eight refined axes (Figure~\ref{fig:axis_heat_corr_pca}D). Together, these visualizations confirm that the refined axes capture stable geometry resolved by subtype. This geometry is primarily characterized by luminality, proliferation, and basal plasticity. Additional functional structure is added by receptor signaling, immunity, and adhesion-invasiveness.

\begin{figure}[!htbp]
    \centering
    \includegraphics[width=1\linewidth]{brca_axis_figures_final/Figure2v2.png}
    \caption{Integrated view of the refined eight-axis representation derived from sample-level axis scores. Each axis score was computed as the mean of gene-wise z-scored expression across the genes present in the corresponding curated biological axis. \textbf{A} Heatmap of mean axis scores by molecular subtype, summarizing the dominant biological programs associated with each class in the reduced axis space. \textbf{B} Spearman correlation matrix among the eight axes, showing the internal dependency structure of the refined representation and the extent of concordance or antagonism between biological programs. \textbf{C} Principal component analysis of the eight-axis space, showing the distribution of individual samples after projection onto the first two principal components and illustrating how subtype structure emerges from this compact biological representation. \textbf{D} One-vs-rest heatmap of subtype-specific effect sizes, expressed as $\Delta$mean (mean axis score in the focal subtype minus the mean score in all remaining samples), highlighting the axes most relatively enriched or depleted in each subtype. Significance symbols indicate Benjamini--Hochberg FDR-adjusted results computed within subtype after two-sided Mann--Whitney U testing (* $q<0.05$, ** $q<0.01$, *** $q<0.001$). Together, these panels show that the curated axis space is internally structured and captures biologically interpretable differences across breast cancer subtypes.}
    \label{fig:axis_heat_corr_pca}
\end{figure}

\subsection{Residual enrichment supported a signaling- and transcription-centered regulatory layer}
An enrichment analysis of the remaining 210 genes in the residual set, after excluding the selected axis assignments, indicated that these genes are associated with receptor and kinase signaling, transcriptional control, protein phosphorylation/modification, and broader intracellular regulatory processes. Representative terms included \emph{Signal Transduction}, \emph{Wnt signaling pathway}, \emph{ErbB signaling pathway}, \emph{MAPK signaling pathway}, \emph{Protein Phosphorylation}, \emph{Signaling by Receptor Tyrosine Kinases}, and several categories related to RNA polymerase II- and DNA-templated transcription (Figure~\ref{fig:enrichment_residual}). These results support the interpretation that the residual layer captures a signaling- and transcription-centered regulatory component rather than unstructured background variation. KEGG terms from this residual set are provided in ST11.



\begin{figure}[!htbp]
        \includegraphics[width=\linewidth]{brca_axis_figures_final/Figure3.png}
    \caption{Functional enrichment analysis of the 210 residual genes retained by the
    GNN-pruned subgraph but not assigned to any predefined biological axis.
    Terms from KEGG, Reactome, and GO Biological Process databases were tested
    using Enrichr (Fisher's exact test, FDR-BH adjusted).
    \textbf{(A)}~Horizontal bar plot showing the top enriched terms ranked by
    statistical significance ($-\log_{10}$ adjusted $p$-value); bar colour
    indicates the annotation source and the annotation at the right of each bar
    reports the number of overlapping genes. The number next to each bar represents the number of genes.
    \textbf{(B)}~Dot plot of the same terms, where the $x$-axis encodes the
    gene ratio (overlap divided by term size), dot diameter is proportional to
    the number of overlapping genes, and colour intensity reflects significance.
    Dominant enrichments include signal transduction (Reactome), Wnt and MAPK
    signalling (KEGG), ErbB signalling, transcriptional regulation, and broad
    cancer pathways, suggesting that the residual gene set captures
    complementary oncogenic programmes beyond the eight curated axes.}
    \label{fig:enrichment_residual}
\end{figure}

Considering the curated axes and the 210 residual genes, we can obtain a structural overview of the nodes associated with the axes (SF4). Therefore, the residual signaling layer is integrated into a broader architecture in which the luminal, receptor-signaling, proliferative, immune, and adhesion-invasion axes are organized.

\newpage
% ============================================================
\section{Discussion}
% ============================================================

We developed an interpretable graph-learning framework for transcriptomic characterization of BC. It combines supervised learning of a sparse, task-relevant subgraph from a comprehensive prior interactome with an axis-based interpretation layer. Its main advantage is that the same learned structure enables both prediction and biological interpretation, unlike existing analyses.

From a problem-solving-protocol perspective, the method was evaluated beyond internal ablations. In addition to full-expression tabular baselines and fixed-prior controls, we evaluated signature-restricted tabular baselines and rewired-prior graph controls within the same cohort-held-out external evaluation setting. This broader comparison strengthens the methodological interpretation in three ways. First, the gains cannot be explained solely by model capacity on the full expression matrix. Several high-capacity tabular models performed well in validation but transferred poorly to the held-out cohorts. For instance, \texttt{XGBoost} dropped from a validation macro-F1 score of 0.975 to a test macro-F1 score of 0.837, whereas the proposed model remained stable (validation 0.928; test 0.924). Second, the gains are not explained only by restricting prediction to a compact retained signature because the strongest 291-gene tabular baseline achieved a test macro-F1 of 0.900, still below the proposed graph-learning model. Third, the gains are not explained by message passing on a fixed or arbitrary graph because fixed-prior and rewired-prior controls remained below the proposed model. Together with the supplementary backbone-replacement diagnostic, these comparisons suggest that supervised topology-aware sparse graph learning has an effect beyond generic deep learning, arbitrary graphs, compact features, and a fixed biological graph.

The ablation results provide relevant methodological insights. The poor performance of \texttt{cond\_only} suggests that sample-conditioned offsets alone are insufficient. A plausible explanation is that instance-level modulation only becomes useful when it operates on top of a stable global ranking of candidate edges. Without typed calibration, heterogeneous interaction classes are subjected to the same selection scale. Thus, sample-specific offsets may increase local selection variability without providing a sufficiently anchored structural prior for robust sparse graph learning. The backbone-replacement diagnostic further supported the use of the ResGAT message-passing block within the FULL InterGATE protocol.


The slight decrease from the validation set to the external test set should not be overinterpreted. In a single-split, cohort-held-out design, a difference of 0.004 macro-F1 is small and likely reflects normal variability in cohort composition and class difficulty rather than significant overfitting. This is because the external test set was never used during supervised training or model selection.

Interestingly, the rewired-prior controls performed similarly to the biological fixed-prior controls, suggesting that prior topology alone is not the main source of performance when graph learning is absent. In other words, running a GNN on a biologically curated graph alone is insufficient. The combination of a biologically constrained candidate prior and supervised sparse topology selection appears to matter more, yielding a task-specific graph that transfers better across held-out cohorts.

A key result is that a compact consensus graph can be extracted from a much larger prior interactome without losing the signal needed for subtype discrimination. Instead of remaining a diffuse regularizer, the retained subgraph becomes an explicit biological entity that can be summarized, visualized, and linked to known BC programs \citep{Barabasi2011NetworkMedicine,Hofree2013NBS,Gogoshin2023GNNCancer}. This is also where our approach differs from conventional expression-based signatures. Clinically established assays, such as PAM50, MammaPrint, and Oncotype~DX, have clear value for subtype assignment, recurrence-risk estimation, and treatment planning in specific patient groups \citep{Parker2009PAM50,VanTVeer2002Signature,Paik2004OncotypeDX}. Our framework is not intended to replace them. Rather, it addresses the need for a topology-aware representation of transcriptomic variation that is both predictive and structurally interpretable. Unlike a fixed gene panel, the present model explicitly learns which parts of the prior interaction network are most relevant to the classification, making the output easier to interpret in terms of regulatory organization than as a simple weighted gene list.

The biological patterns identified through the axis-based interpretation layer align with existing knowledge of BC heterogeneity. The strongest cohort separation followed luminal, proliferative, and basal-associated dimensions, consistent with the long-standing intrinsic subtype framework and later large-scale integrative characterizations of BC \citep{Perou2000MolecularPortraits,Sorlie2003BreastSubtypes,TCGA2012BRCA}. This agreement matters because it suggests that the learned graph is not merely an abstract decision boundary in a high-dimensional feature space. Rather, it retains regulatory components that remain biologically recognizable when projected back into interpretable transcriptional programs.

Several limitations should also be kept in mind. The study is based on the retrospective integration of public transcriptomic cohorts generated on different platforms. Although harmonization and batch correction were applied, residual technical heterogeneity likely remains and may influence model fitting and downstream interpretation. Additionally, the class labels used are meaningful and widely used, but they do not capture the full clinical and molecular complexity of BC. The prior graph is also limited by the completeness and reliability of the interaction resources used to build it. Therefore, the retained subgraph should be understood as a task-specific selection from the available prior knowledge, not as a complete map of BC regulatory biology. From a methodological perspective, performance is also expected to depend on design choices fixed in the primary configuration, such as the prior interactome, sparsity floor \(\kappa\), consensus threshold on selection frequency, and the specific cohort partition used for external transfer evaluation. We addressed some of this sensitivity through predefined architectural ablations, repeated stability-selection runs, and reporting the overlap between runs and consensus. However, broader sensitivity analyses across alternative priors, sparsity levels, external cohort partitions, and repeated alternative-backbone runs remain important future work. Another limitation is the manual curation required for the second-layer axis framework. This was intentional because it enabled a more interpretable and biologically meaningful layer for domain experts. However, any manual curation step introduces analyst judgment and may reflect current pathway knowledge unevenly. For this reason, the refined axes should be viewed as a structured interpretive layer built on top of the learned graph, not as a model-free discovery.

Overall, InterGATE shows that a supervised GNN protocol can turn a large prior interactome into a compact, reproducible, and biologically interpretable subgraph while retaining competitive external performance. Beyond the BC case study, the main contribution is a reproducible framework for learning sparse, task-specific subnetworks from high-dimensional transcriptomic data conditioned by prior biological graphs.


% ============================================================
\section*{Data availability}
% ============================================================

All transcriptomic cohorts used in this study are publicly available. GEO datasets are identified by the accession numbers listed in the main text and supplementary material, and the TCGA BC cohort corresponds to TCGA-BRCA. Data and reproducibility materials required to reproduce the analyses are available from Zenodo under DOI \url{https://doi.org/10.5281/zenodo.20815745}. The final retained-gene universe is provided as Supplementary Data File~1 (SD1).

% ============================================================
\section*{Code availability and implementation}
% ============================================================

Source code for model training, benchmark evaluation, graph export, downstream interpretation, and figure generation is available at \url{https://github.com/alex241278/InterGATE}. The accompanying self-contained package includes \texttt{pyproject.toml}, \texttt{environment.yml}, \texttt{requirements.txt}, \texttt{configs/default.yaml}, a Zenodo download script, a pre-flight setup checker, and an ordered notebook/script workflow.

\begin{comment}  

% ============================================================
\section*{Author biographies}
% ============================================================

Lucia Alvarez-Frutos is a predoctoral researcher at the Instituto de Biomedicina y Genética Molecular, working on cancer biology and computational analysis of transcriptomic data. Laura Senovilla is a researcher at the Instituto de Biomedicina y Genética Molecular, with expertise in cancer biology, tumour heterogeneity, and translational oncology. Alejandro Alvaro-Meca is a public health and biomedical data scientist at Universidad Rey Juan Carlos and CIBERINFEC, with expertise in statistical learning, graph-based modelling, and computational epidemiology.
\end{comment}


% ============================================================
\section*{Funding}
% ============================================================

LS is supported by the Spanish Agencia Estatal de Investigación, AEI (grant number PID2021-126426OB-I00) and the Institut National du Cancer (INCa, PLBIO2022-093). LAF holds a predoctoral fellowship from the Fundación Científica Asociación Española Contra el Cáncer (AECC, PRDVL234237ALVA). 

% ============================================================
\section*{Conflict of Interest}
% ============================================================

The authors declare no conflict of interest.

% ============================================================
\section*{Use of AI-assisted language editing}
% ============================================================

During manuscript preparation, the authors used ChatGPT only as an aid for language editing and readability improvement. No AI system was used to generate the scientific results, figures, tables, references, or final interpretation of the study. All manuscript text was reviewed and edited by the authors, who take full responsibility for the content of the submission.

% ============================================================
%\section*{Supplementary data}
% ============================================================

%Supplementary data are provided separately for review and should be submitted as an independent supplementary file.

%\bibliographystyle{oup-plain}
\bibliographystyle{oup-plain}
\bibliography{Biblio}

\end{document}